\section{Regular Singular Points: Euler Equations}

\noindent
\bigbox{This section sets the stage for Sections~1.5,
1.6,
and 1.7. If you're not interested in those sections, but wish
to learn about Euler equations, omit the introductory paragraphs
and start reading at Definition~\ref{thmtype:7.4.2}.}

\noindent
In the next three sections we'll continue to study  equations of
the form
\begin{equation}\label{eq:7.4.1}
P_0(x)y''+P_1(x)y'+P_2(x)y=0
\end{equation}
where $P_0$, $P_1$, and $P_2$ are polynomials, but the emphasis will
be different from that of Sections~7.2 and 7.3,
where we
obtained solutions of \eqref{eq:7.4.1} near an ordinary point $x_0$ in the
form of power series in $x-x_0$. If $x_0$ is a singular point of
\eqref{eq:7.4.1} (that is, if $P(x_0)=0$),  the solutions can't in
general be represented by power series in $x-x_0$. Nevertheless, it's
often necessary in physical applications to study the behavior of
solutions of \eqref{eq:7.4.1} near a singular point. Although this can be
difficult in the absence of some sort of assumption on the nature of
the singular point, equations that satisfy the requirements of the
next definition can be solved by series methods discussed in the
next three sections. Fortunately, many equations arising in
applications satisfy these requirements.

\begin{definition}\color{blue}\label{thmtype:7.4.1}\rm
Let $P_0$, $P_1$, and $P_2$ be polynomials with no common factor and
suppose  $P_0(x_0)=0$. Then $x_0$ is a {\color{blue}\it regular singular
point\/} of the equation
\begin{equation} \label{eq:7.4.2}
P_0(x)y''+P_1(x)y'+P_2(x)y=0
\end{equation}
if $\eqref{eq:7.4.2}$ can be written as
\begin{equation} \label{eq:7.4.3}
(x-x_0)^2A(x)y''+(x-x_0)B(x)y'+C(x)y=0
\end{equation}
where $A$, $B$, and $C$ are polynomials and $A(x_0)\ne0$;   otherwise,
$x_0$ is an {\color{blue}\it irregular\/} singular point of \eqref{eq:7.4.2}.
\end{definition}

\begin{example}\label{example:7.4.1} \rm
Bessel's equation,
\begin{equation} \label{eq:7.4.4}
x^2y''+xy'+(x^2-\nu^2)y=0,
\end{equation}
has the singular point $x_0=0$. Since this equation is of the form
\eqref{eq:7.4.3} with $x_0=0$, $A(x)=1$, $B(x)=1$, and $C(x)=x^2-\nu^2$,
it follows that $x_0=0$ is a regular singular point of
\eqref{eq:7.4.4}.
 \end{example}

\begin{example}\label{example:7.4.2} \rm
Legendre's equation,
\begin{equation} \label{eq:7.4.5}
(1-x^2)y''-2xy'+\alpha(\alpha+1)y=0,
\end{equation}
has the singular points $x_0=\pm1$. Multiplying through by $1-x$
yields
$$
(x-1)^2(x+1)y''+2x(x-1)y'-\alpha(\alpha+1)(x-1)y=0,
$$
which is of the form \eqref{eq:7.4.3} with $x_0=1$, $A(x)=x+1$,
$B(x)=2x$, and $C(x)=-\alpha(\alpha+1)(x-1)$. Therefore $x_0=1$
is a regular singular point of \eqref{eq:7.4.5}. We leave it to
you to show that $x_0=-1$ is also a regular singular point of
\eqref{eq:7.4.5}.
\end{example}
\newpage

\begin{example}\label{example:7.4.3} \rm
The equation
$$
x^3y''+xy'+y=0
$$
has an irregular  singular point at $x_0=0$. (Verify.)
\end{example}

For convenience we restrict our attention to the case where $x_0=0$ is
a regular singular point of \eqref{eq:7.4.2}. This isn't  really a
restriction, since if $x_0\ne0$ is a regular singular point of
\eqref{eq:7.4.2} then  introducing the new independent variable $t=x-x_0$
and the new unknown $Y(t)=y(t+x_0)$ leads to a differential equation
with polynomial coefficients that has a regular singular point at
$t_0=0$. This is illustrated in Exercise~\ref{exer:7.4.22} for Legendre's
equation, and in Exercise~\ref{exer:7.4.23} for the general case.

\boxit{Euler Equations}

\noindent
The simplest kind of equation with a regular singular point at $x_0=0$
is the Euler equation, defined as follows.

\begin{definition}\color{blue} \label{thmtype:7.4.2} \rm
An
\href{https://mathshistory.st-andrews.ac.uk/Biographies/Euler/}
{\color{blue}\it Euler equation\/} is an equation that can be written in the
form
\begin{equation} \label{eq:7.4.6}
ax^2y''+bxy'+cy=0,
\end{equation}
where $a$, $b$, and $c$ are real constants and $a\ne0$.
\end{definition}

Theorem~\ref{thmtype:5.1.1} implies that \eqref{eq:7.4.6} has
solutions defined on $(0,\infty)$ and $(-\infty,0)$,
since \eqref{eq:7.4.6} can be rewritten as
$$
ay''+{b\over x}y'+{c\over x^2}y=0.
$$
For convenience we'll restrict our attention to the interval
$(0,\infty)$. (Exercise~\ref{exer:7.4.19} deals with solutions of
\eqref{eq:7.4.6} on $(-\infty,0)$.) The key to finding solutions on
$(0,\infty)$ is  that if $x>0$ then $x^r$ is defined as a
real-valued function on $(0,\infty)$ for all values of $r$, and
substituting $y=x^r$ into \eqref{eq:7.4.6} produces
\begin{equation} \label{eq:7.4.7}
\begin{array}{lcl}
ax^2(x^r)''+bx(x^r)'+cx^r&=&ax^2r(r-1)x^{r-2}+bxrx^{r-1}+cx^r\\
&=&[ar(r-1)+br+c]x^r.
\end{array}
\end{equation}
The polynomial
$$
p(r)=ar(r-1)+br+c
$$
is called the {\color{blue}\it indicial polynomial\/} of \eqref{eq:7.4.6}, and
$p(r)=0$
is its {\color{blue}\it indicial equation\/}. From \eqref{eq:7.4.7} we can
see that
$y=x^r$ is a solution of \eqref{eq:7.4.6} on $(0,\infty)$ if and only if
$p(r)=0$. Therefore, if the indicial equation has distinct real roots
$r_1$ and $r_2$ then $y_1=x^{r_1}$ and $y_2=x^{r_2}$ form a
fundamental set of solutions of \eqref{eq:7.4.6} on
$(0,\infty)$, since $y_2/y_1=x^{r_2-r_1}$ is nonconstant.
In this case
$$
y=c_1x^{r_1}+c_2x^{r_2}
$$
is the general solution of \eqref{eq:7.4.6} on $(0,\infty)$.

\begin{example}\label{example:7.4.4} \rm
 Find the general solution of
\begin{equation} \label{eq:7.4.8}
x^2y''-xy'-8y=0
\end{equation}
on  $(0,\infty)$.
\end{example}

\solution
 The indicial polynomial of \eqref{eq:7.4.8} is
$$
p(r)=r(r-1)-r-8=(r-4)(r+2).
$$
Therefore $y_1=x^4$ and $y_2=x^{-2}$ are solutions of \eqref{eq:7.4.8} on
$(0,\infty)$, and its general solution on $(0,\infty)$ is
$$
y=c_1x^4+{c_2\over x^2}.
$$

\begin{example}\label{example:7.4.5} \rm
  Find the general solution of
\begin{equation} \label{eq:7.4.9}
6x^2y''+5xy'-y=0
\end{equation}
on  $(0,\infty)$.
\end{example}

\solution
 The indicial polynomial of \eqref{eq:7.4.9} is
$$
p(r)=6r(r-1)+5r-1=(2r-1)(3r+1).
$$
Therefore the general solution of \eqref{eq:7.4.9} on
$(0,\infty)$ is
$$
y=c_1x^{1/2}+c_2x^{-1/3}. \bbox
$$

If the indicial equation  has a repeated root $r_1$, then
$y_1=x^{r_1}$ is a solution of
\begin{equation} \label{eq:7.4.10}
ax^2y''+bxy'+cy=0,
\end{equation}
on $(0,\infty)$, but \eqref{eq:7.4.10} has no other solution of the form
$y=x^r$. If the indicial equation has complex conjugate zeros then
\eqref{eq:7.4.10} has no real--valued solutions of the form $y=x^r$.
Fortunately we can use the results of Section~5.2 for
constant coefficient equations to solve \eqref{eq:7.4.10} in any case.

\begin{theorem}\color{blue} \label{thmtype:7.4.3}
Suppose the roots of the indicial equation
\begin{equation}
\label{eq:7.4.11} ar(r-1)+br+c=0
\end{equation}
are $r_1$ and $r_2$. Then the general solution of the Euler equation
\begin{equation} \label{eq:7.4.12}
ax^2y''+bxy'+cy=0
\end{equation}
 on  $(0,\infty)$ is
\begin{eqnarray*}
y&=&c_1x^{r_1}+c_2x^{r_2}\mbox{ if $r_1$ and $r_2$ are distinct
real numbers  };
\\ y&=&x^{r_1}(c_1+c_2\ln x)\mbox{ if
$r_1=r_2$  };
\\ y&=&x^{\lambda}\left[c_1\cos\left(\omega\ln x\right)+
c_2\sin\left(\omega\ln x \right)\right]\mbox{ if
$r_1,r_2=\lambda\pm i\omega$ with $\omega>0$}.
\end{eqnarray*}
\end{theorem}

\proof We first show that $y=y(x)$ satisfies \eqref{eq:7.4.12} on
$(0,\infty)$ if and only if $Y(t)=y(e^t)$ satisfies the constant
coefficient equation
\begin{equation} \label{eq:7.4.13}
a{d^2Y\over dt^2}+(b-a){dY\over dt}+cY=0
\end{equation} on
$(-\infty,\infty)$. To do this, it's convenient to write $x=e^t$, or,
equivalently, $t=\ln x$;   thus, $Y(t)=y(x)$, where $x=e^t$. From the
chain rule,
$$
{dY\over dt}={dy\over dx}{dx\over dt}
$$
and, since
$$
{dx\over dt}=e^t=x,
$$
it follows that
\begin{equation} \label{eq:7.4.14}
{dY\over dt}=x{dy\over dx}.
\end{equation}
Differentiating this with
respect to $t$ and using the chain rule again yields
\begin{eqnarray*}
{d^2Y\over dt^2}&=&{d\phantom{t}\over dt}\left(dY\over
dt\right)={d\phantom{t}\over dt}\left(x{dy\over dx}\right)\\[2\jot]
&=&{dx\over dt}{dy\over dx}+x{d^2y\over dx^2}{dx\over dt}\\[2\jot]
&=&x{dy\over dx}+x^2{d^2y\over dx^2}\quad\left(\mbox{ since } {dx\over
dt}=x\right).
\end{eqnarray*}
From this and \eqref{eq:7.4.14},
$$
x^2{d^2y\over dx^2}={d^2Y\over dt^2}-{dY\over dt}.
$$
Substituting this and \eqref{eq:7.4.14} into \eqref{eq:7.4.12} yields
\eqref{eq:7.4.13}. Since \eqref{eq:7.4.11} is the characteristic equation of
\eqref{eq:7.4.13}, Theorem~5.2.1 implies that the general solution
of \eqref{eq:7.4.13} on $(-\infty,\infty)$ is
\begin{eqnarray*}
Y(t)&=&c_1e^{r_1t}+c_2e^{r_2t}\mbox{ if $r_1$ and
$r_2$ are distinct real numbers;  }\\
 Y(t)&=&e^{r_1t}(c_1+c_2t)\mbox{ if
$r_1=r_2$;  }\\
Y(t)&=&e^{\lambda t }\left(c_1\cos\omega
t+c_2\sin\omega t \right)\mbox{ if $r_1,r_2=\lambda\pm i\omega$ with
$\omega\ne0$}.
\end{eqnarray*}
Since $Y(t)=y(e^t)$, substituting $t=\ln x$ in the last three
equations shows that the general solution of \eqref{eq:7.4.12} on
$(0,\infty)$ has the form stated in the theorem.

\begin{example}\label{example:7.4.6} \rm
 Find the general solution of
\begin{equation} \label{eq:7.4.15}
x^2y''-5xy'+9y=0
\end{equation}
on $(0,\infty)$.
\end{example}

\solution
  The indicial polynomial of \eqref{eq:7.4.15} is
$$
p(r)=r(r-1)-5r+9=(r-3)^2.
$$
Therefore the general solution of \eqref{eq:7.4.15} on $(0,\infty)$ is
$$
y=x^3(c_1+c_2 \ln x).
$$

\begin{example}\label{example:7.4.7} \rm
  Find the general solution of
\begin{equation} \label{eq:7.4.16}
x^2y''+3xy'+2y=0
\end{equation}
on  $(0,\infty)$.
\end{example}

\solution
The indicial polynomial of \eqref{eq:7.4.16} is
$$
p(r)=r(r-1)+3r+2=(r+1)^2+1.
$$
The roots of the indicial equation are $r=-1 \pm i$ and the general
solution of \eqref{eq:7.4.16} on $(0,\infty)$ is
$$
y={1\over x}\left[c_1\cos(\ln x)+c_2\sin(\ln x)\right].
$$

\exercises
In Exercises~\ref{exer:7.4.1}--\ref{exer:7.4.18} find the general solution of
the given Euler equation on $(0,\infty)$.

\begin{exerciselist}

 \begin{tabular}[t]{@{}p{168pt}@{}p{168pt}}
\item\label{exer:7.4.1} $x^2y''+7xy'+8y=0$
&\vspace*{10pt}\item\label{exer:7.4.2} $x^2y''-7xy'+7y=0$
\end{tabular}

\begin{tabular}[t]{@{}p{168pt}@{}p{168pt}}
\item\label{exer:7.4.3} $x^2y''-xy'+y=0$
&\item\label{exer:7.4.4} $x^2y''+5xy'+4y=0$
\end{tabular}

\begin{tabular}[t]{@{}p{168pt}@{}p{168pt}}
\item\label{exer:7.4.5} $x^2y''+xy'+y=0$
&\item\label{exer:7.4.6} $x^2y''-3xy'+13y=0$
\end{tabular}

\begin{tabular}[t]{@{}p{168pt}@{}p{168pt}}
\item\label{exer:7.4.7} $x^2y''+3xy'-3y=0$&\item\label{exer:7.4.8}
$12x^2y''-5xy''+6y=0$
\end{tabular}

\begin{tabular}[t]{@{}p{168pt}@{}p{168pt}}
\item\label{exer:7.4.9} $4x^2y''+8xy'+y=0$&\item\label{exer:7.4.10}
$3x^2y''-xy'+y=0$
\end{tabular}

\begin{tabular}[t]{@{}p{168pt}@{}p{168pt}}
\item\label{exer:7.4.11} $2x^2y''-3xy'+2y=0$&\item\label{exer:7.4.12}
$x^2y''+3xy'+5y=0$
\end{tabular}

\begin{tabular}[t]{@{}p{168pt}@{}p{168pt}}
\item\label{exer:7.4.13} $9x^2y''+15xy'+y=0$&
\item\label{exer:7.4.14}
$x^2y''-xy'+10y=0$
\end{tabular}

\begin{tabular}[t]{@{}p{168pt}@{}p{168pt}}
\item\label{exer:7.4.15} $x^2y''-6y=0$&\item\label{exer:7.4.16}
$2x^2y''+3xy'-y=0$
\end{tabular}

\begin{tabular}[t]{@{}p{168pt}@{}p{168pt}}
\item\label{exer:7.4.17} $x^2y''-3xy'+4y=0$&\item\label{exer:7.4.18}
$2x^2y''+10xy'+9y=0$
\end{tabular}

\item\label{exer:7.4.19}
\begin{alist}
\item % (a)
Adapt the proof of Theorem~\ref{thmtype:7.4.3} to show that  $y=y(x)$
satisfies the Euler equation
\setcounter{equation}{0}
\begin{equation} \label{eq:7.4.{exer:7.4.19}A}
ax^2y''+bxy'+cy=0
\end{equation}
on $(-\infty,0)$
if and only if  $Y(t)=y(-e^t)$
$$
a {d^2Y\over dt^2}+(b-a){dY\over dt}+cY=0.
$$
on   $(-\infty,\infty)$.
\item % (b)
Use \part{a}  to show that  the general solution of
\eqref{eq:7.4.{exer:7.4.19}A} on $(-\infty,0)$ is
\begin{eqnarray*}
y&=&c_1|x|^{r_1}+c_2|x|^{r_2}\mbox{ if $r_1$ and $r_2$ are distinct
real numbers;  }
\\ y&=&|x|^{r_1}(c_1+c_2\ln|x|)\mbox{ if
$r_1=r_2$;  }
\\ y&=&|x|^{\lambda}\left[c_1\cos\left(\omega\ln|x|\right)+
c_2\sin\left(\omega\ln|x| \right)\right]\mbox{ if
$r_1,r_2=\lambda\pm i\omega$ with $\omega>0$}.
\end{eqnarray*}
\end{alist}

\item\label{exer:7.4.20} Use reduction of order to show that if
$$
ar(r-1)+br+c=0
$$
has a repeated root $r_1$ then $y=x^{r_1}(c_1+c_2\ln x)$ is the
general solution of
$$
ax^2y''+bxy'+cy=0
$$
on $(0,\infty)$.

\item\label{exer:7.4.21}
  A nontrivial solution of
$$
P_0(x)y''+P_1(x)y'+P_2(x)y=0
$$
is said to be {\color{blue}\it oscillatory} on an interval $(a,b)$ if it has infinitely
many zeros on $(a,b)$.  Otherwise $y$ is said to be {\color{blue}\it nonoscillatory\/}
on $(a,b)$.  Show that the equation
$$
x^2y''+ky=0 \quad (k=\; \mbox{constant})
$$
has oscillatory solutions on $(0,\infty)$ if and only if $k>1/4$.

\item\label{exer:7.4.22}
In Example~\ref{example:7.4.2} we saw  that $x_0=1$ and $x_0=-1$
are regular singular points of Legendre's equation
$$
(1-x^2)y''-2xy'+\alpha(\alpha+1)y=0.
\eqno{\rm (A)}
$$
\begin{alist}
\item % (a)
Introduce the new variables $t=x-1$  and $Y(t)=y(t+1)$, and show that
$y$  is a solution of (A) if and only if
$Y$ is a solution of
$$
t(2+t){d^2Y\over dt^2}+2(1+t){dY\over dt}-\alpha(\alpha+1)Y=0,
$$
which has a regular singular point at $t_0=0$.
\item % (b)
Introduce the new variables $t=x+1$  and $Y(t)=y(t-1)$, and show that
$y$  is a solution of (A) if and only if
$Y$ is a solution of
$$
t(2-t){d^2Y\over dt^2}+2(1-t){dY\over dt}+\alpha(\alpha+1)Y=0,
$$
which has a regular singular point at $t_0=0$.
\end{alist}

\item\label{exer:7.4.23}
Let  $P_0$, $P_1$, and $P_2$ be polynomials  with no common factor,  and
suppose   $x_0\ne0$ is a singular point of
$$
P_0(x)y''+P_1(x)y'+P_2(x)y=0.
\eqno{\rm (A)}
$$
Let $t=x-x_0$ and $Y(t)=y(t+x_0)$.
\begin{alist}
\item % (a)
Show that $y$ is a solution of (A) if and only if $Y$
is a solution of
$$
R_0(t){d^2Y\over dt^2}+R_1(t){dY\over dt}+R_2(t)Y=0.
\eqno{\rm (B)}
$$
where
$$
R_i(t)=P_i(t+x_0),\quad i=0,1,2.
$$
\item % (b)
Show that $R_0$, $R_1$, and $R_2$ are polynomials in $t$ with no common
factors, and $R_0(0)=0$;   thus,  $t_0=0$ is a singular
point of (B).
\end{alist}

\end{exerciselist}
